\begin{table}[h]
\centering
\caption{Figure 3: End-to-end ANCRE pipeline
performance under single-pass streaming conditions.}

\textbf{Measured results:}
\begin{tabular}{lrrl}
\toprule
$n$ & Runtime (s) & Peak RAM & Notes \\
\midrule
$1\times10^2$ & 0.002 & $<0.1$ MB & Python prototype \\
$1\times10^3$ & 0.018 & 0.1 MB    & Python prototype \\
$1\times10^4$ & 0.184 & 1.4 MB    & Python prototype \\
$1\times10^5$ & 1.919 & 12.6 MB   & Python prototype \\
$5\times10^5$ & 9.854 & 59.4 MB   & Python prototype \\
$1\times10^6$ & 0.60  & $<60$ MB  & Rust implementation \\
$3\times10^6$ & 1.80  & $<60$ MB  & Rust implementation \\
\bottomrule
\end{tabular}

\vspace{0.5em}
\textbf{Model-based extrapolation:}
\begin{tabular}{lll}
\toprule
Configuration & Estimated Runtime & Assumptions \\
\midrule
Python @ $3\times10^6$ & $\sim$60s & Linear $O(n)$ projection \\
Rust steady-state & 0.3–1.8s & Warm-cache conditions \\
\bottomrule
\end{tabular}

\vspace{0.5em}
\textbf{Output semantics:}
Single-pass aggregation, constant-size scalar
release (f64), $\varepsilon=0.5$-DP budget,
raw samples destroyed after aggregation,
bounded memory footprint observed experimentally.
\end{table}

\begin{quote}
\textit{The Python implementation is used as a
reference prototype and exhibits linear scaling
dominated by runtime overhead. The Rust
implementation achieves sub-second to low-second
processing for multi-million-sample workloads
under warm-cache conditions while maintaining
bounded memory usage and constant-size DP output.}
\end{quote}
